\documentclass[11pt]{amsart}

\usepackage[T1]{fontenc}
\usepackage{amsmath,amssymb,amsthm}
\usepackage{enumitem}
\usepackage[margin=1.15in]{geometry}
\usepackage[expansion=false]{microtype}
\usepackage{xcolor}
\usepackage[colorlinks=true,linkcolor=blue!55!black,citecolor=blue!55!black]{hyperref}

\newcommand{\NN}{\mathbb{N}}
\newcommand{\CC}{\mathbb{C}}
\newcommand{\RR}{\mathbb{R}}
\newcommand{\cH}{\mathcal{H}}
\newcommand{\cD}{\mathcal{D}}
\newcommand{\cO}{\mathcal{O}}
\newcommand{\Sfin}{S_{\mathrm{fin}}}
\newcommand{\cc}{c_{00}}
\newcommand{\czero}{c_{0}}
\DeclareMathOperator{\Seq}{Seq}
\DeclareMathOperator{\supp}{supp}

\theoremstyle{plain}
\newtheorem{theorem}{Theorem}[section]
\newtheorem{corollary}[theorem]{Corollary}
\newtheorem{proposition}[theorem]{Proposition}
\newtheorem{lemma}[theorem]{Lemma}
\theoremstyle{definition}
\newtheorem{definition}[theorem]{Definition}
\theoremstyle{remark}
\newtheorem{remark}[theorem]{Remark}
\newtheorem{example}[theorem]{Example}

\begin{document}

\title[Rigidity of Symmetric Holomorphic Functions]{Rigidity of Symmetric
Holomorphic Functions on Infinite-Dimensional Sequence Spaces}

\author{David Feldman}
\address{Department of Mathematics and Statistics, University of New
Hampshire, Durham, NH 03824, USA}
\email{david.feldman@unh.edu}

\author{Jon Bannon}
\address{Department of Mathematics, Siena College, Loudonville, NY 12211, USA}
\email{jbannon@siena.edu}

\date{\today}

\subjclass[2020]{Primary 46G20; Secondary 32A05, 46E50, 46B45}
\keywords{Infinite-dimensional holomorphy, symmetric functions,
permutation invariance, rigidity, null sequences, Banach limits,
zero-set map}

\begin{abstract}
We study holomorphic functions on permutation-invariant subsets of an
infinite-dimensional sequence space that are invariant under all finite
permutations of coordinates. Working on domains modeled on the space
$\czero$ of null sequences, we show that permutation symmetry forces the
first differential to vanish along every constant tail: the coordinate
derivatives agree there and an averaging argument annihilates their common
value. This yields rigidity---constancy---on domains of constant sequences,
and more generally on any $\cc$-chain-connected domain on which the
differential vanishes on the finitely supported directions. The method has
a sharp boundary: a precise obstruction prevents the elementary argument
from reaching the finitely many exceptional coordinates of a non-constant
point, so unconditional rigidity on general eventually-constant domains
remains open. In the other direction, an explicit example---any Banach
limit on $\ell^\infty$---shows that the chain-connectedness and boundedness
hypotheses cannot be dropped: mere topological connectedness admits
non-constant symmetric entire functions.
We record the vector-valued corollary and discuss the motivating
application to holomorphic cross-sections of the zero-set map, isolating the
lifting problem that remains open.
\end{abstract}

\maketitle

\section{Introduction}
\label{sec:intro}

In finite dimensions the holomorphic functions on $\CC^n$ invariant under
all permutations of coordinates form a rich algebra, freely generated by
the elementary symmetric polynomials $e_1,\ldots,e_n$. In infinite
dimensions this richness is sharply curtailed: on the natural sequence
spaces built from null sequences, a symmetric holomorphic function is
severely constrained, and on the basic domains is forced to be constant.
The mechanism is that the elementary symmetric functions
$e_k=\sum_{|I|=k}\prod_{i\in I}x_i$ cease to converge for generic
sequences, and this failure is visible already at the level of first
derivatives.

This rigidity phenomenon holds under precise hypotheses, and it has a
definite boundary beyond which it fails. That boundary is real rather than
apparent: the global statement that \emph{every} symmetric holomorphic
function on a connected permutation-invariant subset of $\ell^\infty$ is
constant is false. Any Banach limit furnishes a counterexample
(Example~\ref{ex:banach}). The rigidity is a feature of the null-sequence
geometry, not of symmetry alone.

\subsection*{The main result, informally}
Write $\cc$ for the finitely supported sequences and $\czero$ for the null
sequences, so that $\overline{\cc}=\czero$ in the supremum norm.
Theorem~\ref{thm:rigidity} states: if $X$ is a locally box-open,
permutation-invariant, $\cc$-chain-connected subset of a translate of
$\czero$, and $f:X\to\CC$ is admissibly holomorphic with bounded
differential, invariant under all finite permutations, \emph{and} has
differential vanishing on $\cc$ at every point, then $f$ is constant. The
last hypothesis is essential. The analytic core discharges it
unconditionally at constant points---hence yields outright rigidity on
domains of constant sequences---and along the constant tail of any
eventually-constant point, but \emph{not} at the finitely many exceptional
coordinates of a non-constant point, where a definite obstruction
intervenes (Remark~\ref{rem:obstruction}).

The proof of the vanishing has two elementary movements. First, at any
point $x$ whose coordinates are eventually equal to a constant, the
coordinate derivatives $Df(x)[e_i]$ agree for all indices $i$ in the
constant tail, because the transposition of two tail indices fixes $x$.
Second, averaging $N$ such tail directions produces a perturbation of
supremum norm $1/N\to 0$; a bounded differential then forces the common tail
value to vanish. Constancy then propagates from vanishing derivatives along
finitely supported segments, and $\cc$-chain-connectedness spreads it across
$X$.

\subsection*{Two structural subtleties}
Two features of the setting govern both the reach and the limits of the
argument.

\emph{The base point moves under permutation.} With the coordinate action
$(x\circ\sigma)_n=x_{\sigma(n)}$, the transposition $\tau_{ij}$ satisfies
$(x+\varepsilon e_i)\circ\tau_{ij}=(x\circ\tau_{ij})+\varepsilon e_j$,
which equals $x+\varepsilon e_j$ only when $x_i=x_j$. In general one
obtains
\[
  Df(x)[e_i]=Df(x\circ\tau_{ij})[e_j],
\]
an identity relating the differentials at \emph{two different} base
points. The conclusion ``all coordinate derivatives at $x$ coincide'' is
therefore valid at points fixed by the relevant transpositions---in
particular along the constant tail of an eventually-constant sequence---but
not, by this step alone, at an arbitrary point.

\emph{Finitely supported segments do not connect $\ell^\infty$.} A finite
chain of $\cc$-perturbations alters only finitely many coordinates, so it
can never join $0$ to $(1,1,1,\ldots)$. Vanishing of the differential on
$\cc$ controls $f$ only within a single coset of $\overline{\cc}=\czero$;
across cosets, symmetric holomorphic functions may vary freely, and
Example~\ref{ex:banach} shows the global statement is false on
$\ell^\infty$.

\subsection*{Relation to existing literature}
The principle that infinite symmetry suppresses analytic variation is
familiar in spirit; it is reminiscent of de Finetti-type rigidity in
probability, where exchangeability forces a mixture representation. The
role of the null-sequence geometry, and the Banach-limit obstruction on
$\ell^\infty$, are treated explicitly here.
The requisite background in infinite-dimensional holomorphy is in
Mujica~\cite{Mujica} and Dineen~\cite{Dineen}.

\subsection*{Organization}
Section~\ref{sec:prelim} fixes the framework and the notion of admissible
holomorphicity. Section~\ref{sec:core} proves the analytic core:
tail-derivative agreement, averaging, and vanishing at constant points.
Section~\ref{sec:rigidity-thm} assembles the main theorem and its
vector-valued corollary. Section~\ref{sec:sharpness} gives the
counterexamples showing each hypothesis is necessary.
Section~\ref{sec:application} discusses the zero-set section problem and its
open lifting lemma. Section~\ref{sec:variants} collects further remarks.

\subsection*{Formal verification}
Every numbered mathematical statement of
Sections~\ref{sec:prelim}--\ref{sec:sharpness} has been formally verified in
Lean~4 against Mathlib: the definitions of Section~\ref{sec:prelim}; the
analytic core (Lemmas~\ref{lem:basepoint} and~\ref{lem:cesaro},
Proposition~\ref{prop:tail}, Corollary~\ref{cor:constant},
Proposition~\ref{prop:segment}); the rigidity theorem
(Theorem~\ref{thm:rigidity}) together with the unconditional vanishing
(Proposition~\ref{prop:uncond}); the vector-valued corollary
(Corollary~\ref{cor:nosection}); both counterexamples of
Section~\ref{sec:sharpness}, including the construction of a
finite-permutation-invariant Banach limit; and the exceptional-coordinate
obstruction of Remark~\ref{rem:obstruction}, which appears in the development
as a proved theorem rather than a heuristic. The development compiles without
\texttt{sorry}, and every theorem depends only on the axioms
\texttt{propext}, \texttt{Classical.choice}, and \texttt{Quot.sound}. In one
respect the formal development exceeds the text:
Corollary~\ref{cor:nosection}(ii) is proved for normed targets with no
separation hypothesis, the Hahn--Banach theorem supplying the separating
functionals. The conditional Proposition~\ref{prop:conditional} is not
formalized as such; its rigidity input is Corollary~\ref{cor:nosection}, and
the Lifting Lemma on which it rests is open.

The formalization was carried out with Aristotle (Harmonic) from written work
orders, archived alongside the proofs. It repaid the effort: an earlier
draft's permutation identity omitted the movement of the base point, and its
local-to-global step was false outright, the development supplying an
explicit counterexample; the statements above reflect those corrections.

The formalization, the work orders, and the audit reports are archived at
Zenodo and developed on GitHub:
\begin{center}
\url{https://doi.org/10.5281/zenodo.21944129}\\
\url{https://github.com/DavidVFeldman/rigidity-symmetric-holomorphic}
\end{center}

\section{Preliminaries}
\label{sec:prelim}

Let $\ell^\infty=\ell^\infty(\NN)$ be the Banach space of bounded complex
sequences with the supremum norm $\|\cdot\|_\infty$. Let $\czero\subset
\ell^\infty$ be the closed subspace of sequences tending to $0$, and let
$\cc\subset\czero$ be the dense subspace of finitely supported sequences.
Write $e_k$ for the sequence with $1$ in position $k$ and $0$ elsewhere,
and recall $\overline{\cc}=\czero$ in $\|\cdot\|_\infty$.

Let $\Sfin$ denote the group of \emph{finite permutations} of $\NN$: the
bijections $\sigma:\NN\to\NN$ fixing all but finitely many elements. It
acts on $\CC^\NN$ by $(x\circ\sigma)_n=x_{\sigma(n)}$.

\begin{definition}
\label{def:domain}
A subset $X\subset\CC^\NN$ is:
\begin{itemize}[leftmargin=2em]
  \item \emph{permutation-invariant} if $x\circ\sigma\in X$ for every
        $x\in X$ and every $\sigma\in\Sfin$;
  \item \emph{locally box-open} if for every $x\in X$ there is $r>0$ with
        $x+h\in X$ for all $h\in\cc$ with $\|h\|_\infty<r$;
  \item \emph{$\cc$-chain-connected} if for any $x,y\in X$ there is a finite
        sequence $x=z_0,z_1,\ldots,z_m=y$ in $X$ with $z_{\ell+1}-z_\ell\in
        \cc$ and the segment $\{z_\ell+t(z_{\ell+1}-z_\ell):t\in[0,1]\}$
        contained in $X$ for each $\ell$.
\end{itemize}
\end{definition}

The distinction between $\cc$-chain-connectedness and topological
connectedness is central; see Section~\ref{sec:sharpness}. For orientation:
$\czero$ itself is $\cc$-chain-connected (any null sequence is a
sup-norm limit of finitely supported ones, but more to the point any two
finitely supported sequences are joined by a single $\cc$-segment, and the
general case follows by density arguments on the relevant domains), whereas
$\ell^\infty$ is not.

\begin{definition}[Admissible holomorphicity]
\label{def:holomorphic}
A function $f:X\to\CC$ on a locally box-open, permutation-invariant set
$X\subset\CC^\NN$ is \emph{admissibly holomorphic} if it carries a
differential datum $x\mapsto Df(x)$, where each $Df(x):\cc\to\CC$ is
$\CC$-linear, subject to:
\begin{enumerate}[label=\textup{(\roman*)}, leftmargin=2.8em]
\item \emph{(Coordinatewise holomorphy.)} For every $x\in X$ and every
      $k\in\NN$, the map $\varepsilon\mapsto f(x+\varepsilon e_k)$ is
      holomorphic in a neighborhood of $0\in\CC$.
\item \emph{(Bounded first-order expansion along $\cc$.)} For every $x\in
      X$ there is a constant $C_x$ with $|Df(x)[h]|\le C_x\|h\|_\infty$ for
      all $h\in\cc$, and
      \[
        f(x+h)=f(x)+Df(x)[h]+o(\|h\|_\infty)
        \qquad\text{as }h\to 0\text{ in }\cc.
      \]
\end{enumerate}
\end{definition}

\begin{remark}[Standard holomorphy is admissible]
\label{rem:standard}
Both standard notions of holomorphy on $\ell^\infty$ (or on $\czero$)
supply Definition~\ref{def:holomorphic} with a \emph{bounded} differential.
A Fr\'echet-holomorphic function has a continuous Fr\'echet differential
$Df(x)$, whose restriction to $\cc$ is bounded and supplies~(ii), with~(i)
immediate. A G\^ateaux-holomorphic and locally bounded function---the
standard definition in infinite dimensions, see~\cite{Mujica,Dineen}---
yields, via the one-variable Cauchy estimates applied coordinatewise, a
bounded linear differential on finitely supported directions, giving both
conditions. The boundedness in~(ii) is thus automatic in every standard
setting; we include it explicitly because it is exactly what the averaging
step requires, and because Definition~\ref{def:holomorphic} without it does
not force sup-norm continuity of $f$ along $\cc$.
\end{remark}

\section{The Analytic Core}
\label{sec:core}

Throughout this section $X\subset\CC^\NN$ is locally box-open and
permutation-invariant, and $f:X\to\CC$ is admissibly holomorphic and
$\Sfin$-invariant.

\subsection{The permutation identity}

\begin{lemma}[Base-point identity]
\label{lem:basepoint}
For distinct $i,j\in\NN$ and every $x\in X$,
\[
  Df(x)[e_i] = Df(x\circ\tau_{ij})[e_j],
\]
where $\tau_{ij}\in\Sfin$ is the transposition of $i$ and $j$. In
particular, if $x_i=x_j$ then $x\circ\tau_{ij}=x$ and
$Df(x)[e_i]=Df(x)[e_j]$.
\end{lemma}

\begin{proof}
For $\varepsilon\in\CC$ with $|\varepsilon|$ small, local box-openness
gives $x+\varepsilon e_i\in X$. A direct computation from
$(y\circ\tau_{ij})_n=y_{\tau_{ij}(n)}$ shows
\[
  (x+\varepsilon e_i)\circ\tau_{ij}
  = (x\circ\tau_{ij}) + \varepsilon e_j .
\]
By $\Sfin$-invariance of $f$,
\[
  f(x+\varepsilon e_i)
  = f\bigl((x+\varepsilon e_i)\circ\tau_{ij}\bigr)
  = f\bigl((x\circ\tau_{ij})+\varepsilon e_j\bigr).
\]
Subtract $f(x)=f(x\circ\tau_{ij})$ (again by invariance), divide by
$\varepsilon$, and let $\varepsilon\to0$; coordinatewise holomorphy,
condition~(i), guarantees the two one-variable limits exist and equal the
respective coordinate derivatives, giving
$Df(x)[e_i]=Df(x\circ\tau_{ij})[e_j]$. If $x_i=x_j$ then swapping positions
$i,j$ leaves $x$ unchanged, so $x\circ\tau_{ij}=x$ and the two base points
coincide.
\end{proof}

\begin{remark}
The identity of Lemma~\ref{lem:basepoint} relates \emph{different} base
points whenever $x_i\ne x_j$, and one cannot conclude
$Df(x)[e_i]=Df(x)[e_j]$ without further input. Equality of coordinate
derivatives at a single point is available exactly along coordinates on
which $x$ is constant.
\end{remark}

\subsection{Averaging}

\begin{lemma}[Ces\`aro vanishing]
\label{lem:cesaro}
Fix $x\in X$ and let $C_x$ bound $Df(x)$ as in
Definition~\ref{def:holomorphic}(ii). For any injective sequence
$i_1,i_2,\ldots$ of indices,
\[
  \frac{1}{N}\sum_{k=1}^N Df(x)[e_{i_k}] \longrightarrow 0
  \qquad (N\to\infty).
\]
\end{lemma}

\begin{proof}
Set $h_N=\tfrac1N\sum_{k=1}^N e_{i_k}\in\cc$. Distinct indices give
$\|h_N\|_\infty=1/N$. By linearity, $\tfrac1N\sum_{k=1}^N
Df(x)[e_{i_k}]=Df(x)[h_N]$, and boundedness gives
$|Df(x)[h_N]|\le C_x\|h_N\|_\infty=C_x/N\to0$.
\end{proof}

\subsection{Vanishing at eventually-constant points}

Say $x\in\CC^\NN$ is \emph{eventually constant with value $c$} if
$x_n=c$ for all but finitely many $n$; write $\supp_c(x)=\{n:x_n\ne c\}$,
a finite set, for its \emph{exceptional set}.

\begin{proposition}[Tail vanishing]
\label{prop:tail}
Let $x\in X$ be eventually constant with value $c$. Then $Df(x)[e_i]=0$ for
every index $i\notin\supp_c(x)$; that is, the coordinate derivative
vanishes along the constant tail.
\end{proposition}

\begin{proof}
Let $T=\NN\setminus\supp_c(x)$ be the (cofinite) tail, so $x_i=c$ for all
$i\in T$. For $i,j\in T$ we have $x_i=x_j=c$, so
Lemma~\ref{lem:basepoint} gives $Df(x)[e_i]=Df(x)[e_j]$. Hence there is a
scalar $\lambda$ with $Df(x)[e_i]=\lambda$ for all $i\in T$. Since $T$ is
infinite, choose an injective sequence $i_1,i_2,\ldots$ inside $T$; each
term of the Ces\`aro average in Lemma~\ref{lem:cesaro} equals $\lambda$, so
the average is constantly $\lambda$ while tending to $0$. Therefore
$\lambda=0$, i.e. $Df(x)[e_i]=0$ for every $i\in T$.
\end{proof}

\begin{corollary}[Vanishing at constant points]
\label{cor:constant}
If the constant sequence $x\equiv c$ lies in $X$, then $Df(x)[h]=0$ for all
$h\in\cc$.
\end{corollary}

\begin{proof}
A constant sequence has empty exceptional set, so
Proposition~\ref{prop:tail} gives $Df(x)[e_k]=0$ for every $k$. Any
$h\in\cc$ is a finite combination $h=\sum_k h_k e_k$, and linearity of
$Df(x)$ gives $Df(x)[h]=0$.
\end{proof}

\subsection{Constancy along finitely supported segments}

\begin{proposition}[Segment rigidity]
\label{prop:segment}
Let $x\in X$ and $h\in\cc$ be such that the segment
$\gamma(t)=x+t\,h$ lies in $X$ for all $t\in[0,1]$, and suppose
$Df(\gamma(t))[h]=0$ for every $t\in[0,1]$. Then $f(x+h)=f(x)$.
\end{proposition}

\begin{proof}
Consider $\varphi(t)=f(x+t\,h)$ on $[0,1]$. Fix $t_0\in[0,1]$. For small
real $s$, the first-order expansion at the base point $\gamma(t_0)$,
applied to the finitely supported increment $s\,h$, gives
\[
  \varphi(t_0+s)-\varphi(t_0)
  = f\bigl(\gamma(t_0)+s\,h\bigr)-f(\gamma(t_0))
  = s\,Df(\gamma(t_0))[h] + o(|s|\,\|h\|_\infty)
  = o(|s|),
\]
using $Df(\gamma(t_0))[h]=0$. Hence $\varphi$ is differentiable at $t_0$
with $\varphi'(t_0)=0$. As $t_0\in[0,1]$ was arbitrary, $\varphi$ has
identically zero derivative on $[0,1]$ and is therefore constant; in
particular $\varphi(1)=\varphi(0)$, i.e. $f(x+h)=f(x)$.
\end{proof}

\section{The Rigidity Theorem}
\label{sec:rigidity-thm}

The natural domains for the main theorem are those sitting inside a single
translate of $\czero$---equivalently, those whose points share one eventual
value. On such a domain the $\cc$-directional derivative structure controls
the function, and the theorem below gives constancy once the differential
vanishes on $\cc$ throughout. Proposition~\ref{prop:uncond} then records the
extent to which that hypothesis holds without further assumption.

\begin{theorem}[Rigidity on null-sequence domains]
\label{thm:rigidity}
Fix $c\in\CC$ and let $X\subset (c\cdot\mathbf 1)+\czero$ be nonempty,
locally box-open, permutation-invariant, and $\cc$-chain-connected, where
$\mathbf 1=(1,1,\ldots)$. Let $f:X\to\CC$ be admissibly holomorphic with
bounded differential and invariant under all finite permutations. Suppose
moreover that $Df(x)[h]=0$ for all $x\in X$ and all $h\in\cc$. Then $f$ is
constant on $X$.
\end{theorem}

\begin{proof}
Given $x,y\in X$, take a $\cc$-chain $x=z_0,\ldots,z_m=y$ as in
Definition~\ref{def:domain}, with each increment
$h_\ell:=z_{\ell+1}-z_\ell\in\cc$ and each segment inside $X$. Along the
$\ell$-th segment the standing hypothesis gives
$Df(z_\ell+t\,h_\ell)[h_\ell]=0$ for all $t\in[0,1]$, so
Proposition~\ref{prop:segment} yields $f(z_{\ell+1})=f(z_\ell)$. Chaining
the equalities gives $f(y)=f(x)$. As $x,y$ were arbitrary, $f$ is constant.
\end{proof}

The theorem carries the hypothesis $Df|_{\cc}\equiv0$ because it is not
implied, at every point of an eventually-constant domain, by the remaining
hypotheses. What symmetry alone yields is the value of the differential at
constant points and along constant tails.

\begin{proposition}[Vanishing at constant points and tails]
\label{prop:uncond}
In the setting of Section~\ref{sec:core}, $Df(c\cdot\mathbf 1)[h]=0$ for
every $h\in\cc$ whenever $c\cdot\mathbf 1\in X$
\textup{(Corollary~\ref{cor:constant})}; and for any eventually-constant
$x\in X$ with value $c$, $Df(x)[e_i]=0$ for every tail index $i$ with
$x_i=c$ \textup{(Proposition~\ref{prop:tail})}.
\end{proposition}

Proposition~\ref{prop:uncond} marks how far the hypothesis
$Df|_{\cc}\equiv0$ of Theorem~\ref{thm:rigidity} holds without further
assumption. In particular, on a domain of constant sequences every
coordinate is a tail coordinate, so the hypothesis holds throughout and $f$
is constant unconditionally. At the exceptional coordinates of a
non-constant point the differential is \emph{not} controlled by this
argument, for a definite reason.

\begin{remark}[The exceptional-coordinate obstruction]
\label{rem:obstruction}
Removing an exceptional coordinate $k$ of a point $x$ (where $x_k\ne c$) by
passing to the ``zeroed'' point $p$ obtained from $x$ by resetting the
$k$-th entry to $c$---at which $k$ is a tail coordinate---does not transport
the vanishing back to $x$. The base-point identity Lemma~\ref{lem:basepoint}
gives $Df(x)[e_k]=Df(x\circ\tau_{kj})[e_j]$, but $x\circ\tau_{kj}$ merely
\emph{moves} the exceptional value from position $k$ to position $j$; it is
not a point at which $k$ is a tail coordinate. Concretely, $x\circ\tau_{kj}$
and $p\circ\tau_{kj}$ disagree at position $j$ (values $x_k$ versus $c$), so
the identity at $x$ does not relate to $p$. Since each eventually-constant
sequence has only finitely many exceptional coordinates, the averaging
argument of Lemma~\ref{lem:cesaro}---which requires infinitely many equal
coordinate derivatives at a \emph{single} base point---does not apply to the
exceptional directions. Whether $Df|_{\cc}\equiv0$ holds at every point
under a stronger analytic hypothesis---for instance continuity of
$x\mapsto Df(x)$ in the supremum norm, available for Fr\'echet-holomorphic
$f$---is open.
\end{remark}

\subsection{The vector-valued corollary}

\begin{corollary}
\label{cor:nosection}
Let $X$ be as in Theorem~\ref{thm:rigidity}, and let $Y$ be a locally
convex complex topological vector space. Suppose $\Sigma:X\to Y$ is
$\Sfin$-equivariant, $\Sigma(x\circ\sigma)=\Sigma(x)$, and such that for
every $\lambda\in Y^*$ the scalar function $\lambda\circ\Sigma$ satisfies
the hypotheses of Theorem~\ref{thm:rigidity}. Then:
\begin{enumerate}[label=\textup{(\roman*)}, leftmargin=2.8em]
\item $\lambda\circ\Sigma$ is constant for every $\lambda\in Y^*$;
\item if $Y^*$ separates the points of $Y$, then $\Sigma$ is constant.
\end{enumerate}
\end{corollary}

\begin{proof}
Each $\lambda\circ\Sigma$ is $\Sfin$-invariant because $\Sigma$ is
equivariant, so Theorem~\ref{thm:rigidity} gives~(i). For~(ii), if
$\Sigma(x)\ne\Sigma(y)$ some $\lambda\in Y^*$ separates the values,
contradicting~(i).
\end{proof}

\section{Sharpness: the Hypotheses Are Necessary}
\label{sec:sharpness}

The hypothesis confining the domain to a single $\czero$-coset, which is to
say the hypothesis of $\cc$-chain-connectedness, cannot be dropped, as the
following two examples show. Both are genuinely symmetric and genuinely holomorphic; what fails is
precisely the passage from $\cc$-local to global constancy.

\begin{example}[Two cosets: a disconnected witness]
\label{ex:twocosets}
For $c\in\{0,1\}$ let $S_c=\{x\in\ell^\infty : x_n=c \text{ for all but
finitely many } n\}$ be the eventually-$c$ sequences, and put $X=S_0\cup
S_1$. Each $S_c$ is permutation-invariant and locally box-open, and a
$\cc$-perturbation of a point of $S_c$ remains in $S_c$; thus $X$ is
locally box-open and permutation-invariant. Define $f:X\to\CC$ by $f\equiv
0$ on $S_0$ and $f\equiv 1$ on $S_1$. Then $f$ is $\Sfin$-invariant, and it
is admissibly holomorphic with $Df\equiv 0$ (it is locally constant along
$\cc$, since $\cc$-perturbations do not change the eventual value). Yet $f$
is not constant. The obstruction is exactly that $S_0$ and $S_1$ are
distinct cosets of $\czero$: no finite chain of $\cc$-steps joins them, so
$X$ is not $\cc$-chain-connected.
\end{example}

\begin{example}[A connected witness: Banach limits]
\label{ex:banach}
Let $L\in(\ell^\infty)^*$ be a Banach limit: a continuous linear functional
with $L(\mathbf 1)=1$, invariant under the shift, and hence under every
finite permutation, and extending the ordinary limit on convergent
sequences. Being continuous and linear, $L$ is entire on $\ell^\infty$, so
it is admissibly holomorphic with $DL(x)=L$ for all $x$; its differential is
bounded (by $\|L\|$). Since $L$ vanishes on $\czero\supset\cc$, we have
$DL(x)|_{\cc}=0$ at every point---exactly as the analytic core predicts.
Nonetheless $L$ is \emph{not} constant: $L(0)=0$ while $L(\mathbf 1)=1$.

Here the domain is all of $\ell^\infty$, which \emph{is} connected (indeed
convex), so mere topological connectedness does not suffice for rigidity.
What fails is $\cc$-chain-connectedness: $0$ and $\mathbf 1$ are at
sup-distance $1$ from one another across the quotient
$\ell^\infty/\czero$, and no finite $\cc$-chain bridges them. This example
shows that the restriction to null-sequence domains in
Theorem~\ref{thm:rigidity} is not a technical convenience but a
necessity.
\end{example}

\begin{remark}[Sharpness of the topology]
\label{rem:ell1}
A complementary degeneration occurs if the supremum norm is replaced by a
summing norm. On $\ell^1\cap D^\NN$ the averages
$h_N=\tfrac1N\sum_{k=1}^N e_{i_k}$ have $\|h_N\|_1=1$ for all $N$, so the
Ces\`aro step (Lemma~\ref{lem:cesaro}) fails, and indeed $f(x)=\sum_k x_k$
is a non-constant symmetric holomorphic function. Thus a topology in which
averaged coordinate perturbations become small---such as the supremum
norm on $\czero$---is essential to the averaging mechanism.
\end{remark}

\section{Application: the Zero-Set Section Problem}
\label{sec:application}

We indicate how the rigidity theorem bears on a natural geometric question,
and we state precisely the point at which the argument becomes conditional.

Let $\cH=\cO(B(0,1))$ be the Fr\'echet space of holomorphic functions on
the open unit disk with the compact-open topology, and let $\cD$ be the
space of admissible infinite discrete subsets of $B(0,1)$---sequences with
no limit point inside the disk. By the Weierstrass factorization theorem
the zero-set map
\[
  \pi:\cH\setminus\{0\}\longrightarrow\cD, \qquad f\longmapsto
  (\text{zero multiset of }f),
\]
is surjective. One asks whether $\pi$ admits a holomorphic cross-section
$\sigma:\cD\to\cH$ with $\pi\circ\sigma=\mathrm{id}$.

Model $\cD$ as a quotient $\Seq/S_\infty$, where $\Seq$ is the space of
locally finite sequences in $B(0,1)$; the naturally arising domain of
sequences is $\cc$-chain-connected, since admissible sets are altered one
point at a time by finitely supported moves. The rigidity theorem then
yields a conditional obstruction.

\begin{proposition}[Conditional obstruction]
\label{prop:conditional}
Suppose a cross-section $\sigma:\cD\to\cH$ lifts to an
$\Sfin$-equivariant map $\widetilde\sigma:X\to\cH$ on the associated
$\cc$-chain-connected sequence domain $X$, with $\lambda\circ\widetilde
\sigma$ admissibly holomorphic (bounded differential, differential
vanishing on $\cc$) for every evaluation functional $\lambda\in\cH^*$.
Then no such $\sigma$ exists.
\end{proposition}

\begin{proof}
Since $\cH$ separates points via evaluations,
Corollary~\ref{cor:nosection}(ii) forces $\widetilde\sigma$ to be constant,
say $\widetilde\sigma\equiv f_0$. But then $\pi(f_0)=\pi(\widetilde\sigma
(x))$ would equal the varying zero set of $x$ for every $x\in X$, which is
absurd.
\end{proof}

The hypothesis is a genuine lifting condition, stated precisely below.

\begin{quotation}\noindent
\textbf{Open Problem (Lifting Lemma).}
Let $Y$ be a locally convex space. Under a suitable complex structure on
$\cD=\Seq/S_\infty$, does every holomorphic map $F:\cD\to Y$ lift to an
$\Sfin$-equivariant map $\widetilde F:\Seq\to Y$ whose scalarizations are
admissibly holomorphic with differentials vanishing on $\cc$?
\end{quotation}

\noindent An affirmative answer would, through
Corollary~\ref{cor:nosection}, upgrade
Proposition~\ref{prop:conditional} to the unconditional statement that
$\pi$ admits no holomorphic section on all of $\cD$. We leave this as the
natural next target.

\begin{remark}[A complementary unconditional obstruction]
Independently of the lifting problem, there is no section of the special
\emph{product} form $\sigma(A)=\prod_{a\in A}E(z,a)$ for a fixed
holomorphic factor $E$: symmetry forces the factor to be uniform, a uniform
factor imposes one fixed summability condition on the zero set for
convergence, and $\cD$ contains admissible sets violating any such
condition. Thus symmetric product sections fail for convergence reasons
unconditionally, while any symmetric non-product section would be obstructed
by rigidity, conditionally on the lifting.
\end{remark}

\section{Further Remarks}
\label{sec:variants}

\begin{enumerate}[label=\textup{(\arabic*)}, leftmargin=2.4em, itemsep=0.5em]

\item \emph{What symmetry alone gives.} The analytic core
(Section~\ref{sec:core}) is valid on any locally box-open
permutation-invariant domain, without chain-connectedness: it yields
tail-derivative vanishing at every point and full vanishing at constant
points. Chain-connectedness enters only to globalize, and
Example~\ref{ex:banach} shows it cannot be omitted.

\item \emph{Relation to de Finetti's theorem.} The result is an analytic
cousin of de Finetti's theorem: exchangeability of a measure on
$\{0,1\}^\NN$ forces dependence only on ``collective'' statistics, whereas
here symmetry forces a holomorphic function on a null-sequence domain to
depend on nothing at all. The Banach-limit example marks the boundary: on
$\ell^\infty$, symmetric ``limit-like'' functionals persist, just as tail
$\sigma$-fields carry exchangeable information the mixture representation
cannot see.

\item \emph{Comparison with finite dimensions.} For each finite $n$ the
$S_n$-invariant holomorphic functions on $\CC^n$ form the polynomial
algebra on $e_1,\ldots,e_n$, each a finite and hence entire sum. In
infinite dimensions these sums diverge on generic sequences;
Proposition~\ref{prop:tail} is the precise derivative-level expression of
that divergence along the constant tail.

\item \emph{Positive results on smaller domains.} On $\ell^2\cap D^\NN$ the
power sums $p_k=\sum_n a_n^k$ converge and are holomorphic for $k\ge3$,
providing non-constant symmetric holomorphic functions. Rigidity is a
feature of the $\czero$ geometry with its averaging mechanism, not of
symmetry as such.

\item \emph{A topological analogue.} The averaging argument is not specific
to holomorphy. The same reasoning shows that a permutation-invariant
function on a $\cc$-chain-connected domain in $\czero$ that is Lipschitz and
G\^ateaux-differentiable along $\cc$, with bounded differential, is
constant. Holomorphy is used only to produce the coordinate derivatives in
Lemma~\ref{lem:basepoint}; the heart of the matter is the metric fact
$\|h_N\|_\infty=1/N\to0$ together with the coset structure of $\czero$
inside $\ell^\infty$.

\end{enumerate}

\begin{thebibliography}{9}

\bibitem{Dineen}
S.~Dineen,
\emph{Complex Analysis on Infinite Dimensional Spaces},
Springer Monographs in Mathematics, Springer-Verlag, London, 1999.

\bibitem{Mujica}
J.~Mujica,
\emph{Complex Analysis in Banach Spaces},
North-Holland Mathematics Studies, vol.~120, North-Holland, Amsterdam, 1986.

\bibitem{Macdonald}
I.~G.~Macdonald,
\emph{Symmetric Functions and Hall Polynomials},
2nd ed., Oxford University Press, Oxford, 1995.

\bibitem{Arnold}
V.~I.~Arnol'd,
On some topological invariants of algebraic functions,
\emph{Trans. Moscow Math. Soc.} \textbf{21} (1970), 30--52.

\end{thebibliography}

\end{document}
